\documentclass[11pt]{article}

\usepackage[margin=0.78in]{geometry}
\usepackage[T1]{fontenc}
\usepackage[utf8]{inputenc}
\usepackage{amsmath}
\usepackage{microtype}
\usepackage{enumitem}
\usepackage[hidelinks]{hyperref}

% Generated from results/paper_stats.json by the paper-statistics builder.
% Generated from results/paper_stats.json; do not edit.
\newcommand{\DiagnosticSlices}{87}
\newcommand{\DiagnosticFlagged}{14}
\newcommand{\DiagnosticInconclusive}{9}
\newcommand{\DiagnosticRows}{87,000}
\newcommand{\DiagnosticThreshold}{0.50}
\newcommand{\RewardModelName}{xlm-roberta-base}
\newcommand{\RewardTrainRows}{50,000}
\newcommand{\RewardDevRows}{5,000}
\newcommand{\RewardSteps}{500}
\newcommand{\RewardCheckpointSeeds}{1}
\newcommand{\PrimaryPool}{18}
\newcommand{\PrimarySegments}{410}
\newcommand{\ExternalPool}{30}
\newcommand{\ExternalSegments}{541}
\newcommand{\PoolPermutations}{100}
\newcommand{\BanditContexts}{541}
\newcommand{\BanditActions}{30}
\newcommand{\BanditTrainContexts}{367}
\newcommand{\BanditTestContexts}{174}
\newcommand{\BanditTestDocuments}{56}
\newcommand{\BanditUpdates}{200}
\newcommand{\BanditSeeds}{5}
\newcommand{\UniformCesa}{66.27}
\newcommand{\WMTPrimaryName}{WMT25}
\newcommand{\WMTExternalName}{WMT26}
\newcommand{\AllExcludedSlices}{0}
\newcommand{\AllRemovedRows}{0}
\newcommand{\AllDevSpearman}{0.438}
\newcommand{\AllPrimaryFull}{-7.86}
\newcommand{\AllPrimaryChange}{-0.58}
\newcommand{\AllExternalFull}{44.70}
\newcommand{\AllExternalChange}{-19.66}
\newcommand{\AllBanditCesa}{33.07}
\newcommand{\AllBanditChange}{-33.19}
\newcommand{\AllBanditLow}{-43.39}
\newcommand{\AllBanditHigh}{-21.39}
\newcommand{\AllTopProbability}{0.734}
\newcommand{\FilteredExcludedSlices}{14}
\newcommand{\FilteredRemovedRows}{239,364}
\newcommand{\FilteredDevSpearman}{0.449}
\newcommand{\FilteredPrimaryFull}{-7.47}
\newcommand{\FilteredPrimaryChange}{-0.20}
\newcommand{\FilteredExternalFull}{46.37}
\newcommand{\FilteredExternalChange}{-17.99}
\newcommand{\FilteredBanditCesa}{33.86}
\newcommand{\FilteredBanditChange}{-32.40}
\newcommand{\FilteredBanditLow}{-43.44}
\newcommand{\FilteredBanditHigh}{-19.88}
\newcommand{\FilteredTopProbability}{0.695}
\newcommand{\RandomExcludedSlices}{14}
\newcommand{\RandomRemovedRows}{235,497}
\newcommand{\RandomDevSpearman}{0.432}
\newcommand{\RandomPrimaryFull}{-7.79}
\newcommand{\RandomPrimaryChange}{-0.52}
\newcommand{\RandomExternalFull}{46.43}
\newcommand{\RandomExternalChange}{-17.93}
\newcommand{\RandomBanditCesa}{39.67}
\newcommand{\RandomBanditChange}{-26.60}
\newcommand{\RandomBanditLow}{-36.34}
\newcommand{\RandomBanditHigh}{-15.22}
\newcommand{\RandomTopProbability}{0.564}
\newcommand{\DirectionExcludedSlices}{14}
\newcommand{\DirectionRemovedRows}{247,332}
\newcommand{\DirectionDevSpearman}{0.434}
\newcommand{\DirectionPrimaryFull}{-7.61}
\newcommand{\DirectionPrimaryChange}{-0.33}
\newcommand{\DirectionExternalFull}{48.04}
\newcommand{\DirectionExternalChange}{-16.32}
\newcommand{\DirectionBanditCesa}{67.92}
\newcommand{\DirectionBanditChange}{+1.66}
\newcommand{\DirectionBanditLow}{+1.23}
\newcommand{\DirectionBanditHigh}{+2.10}
\newcommand{\DirectionTopProbability}{0.056}
\newcommand{\FilteredMinusAllPrimary}{+0.39}
\newcommand{\FilteredMinusAllPrimaryLow}{-0.41}
\newcommand{\FilteredMinusAllPrimaryHigh}{+1.25}
\newcommand{\FilteredMinusAllExternal}{+1.67}
\newcommand{\FilteredMinusAllExternalLow}{+0.01}
\newcommand{\FilteredMinusAllExternalHigh}{+3.42}
\newcommand{\FilteredMinusRandomPrimary}{+0.32}
\newcommand{\FilteredMinusRandomPrimaryLow}{-0.43}
\newcommand{\FilteredMinusRandomPrimaryHigh}{+1.10}
\newcommand{\FilteredMinusRandomExternal}{-0.06}
\newcommand{\FilteredMinusRandomExternalLow}{-1.86}
\newcommand{\FilteredMinusRandomExternalHigh}{+1.75}
\newcommand{\FilteredMinusDirectionPrimary}{+0.14}
\newcommand{\FilteredMinusDirectionPrimaryLow}{-0.54}
\newcommand{\FilteredMinusDirectionPrimaryHigh}{+0.85}
\newcommand{\FilteredMinusDirectionExternal}{-1.67}
\newcommand{\FilteredMinusDirectionExternalLow}{-3.69}
\newcommand{\FilteredMinusDirectionExternalHigh}{+0.32}


\setlength{\parindent}{0pt}
\setlength{\parskip}{0.45em}
\setlist[itemize]{leftmargin=1.35em,itemsep=0.15em,topsep=0.2em}
\pagestyle{plain}

\newcommand{\field}[1]{\vspace{0.25em}\textbf{#1}\quad}

\title{Can We Predict When Optimizing a Translation Reward Will Hurt?}
\author{Isaac Thompson\\Reinforcement Learning Project Proposal}
\date{\today}

\begin{document}
\maketitle
\vspace{-1.7em}

\field{Topic:}
Reward modeling, proxy objectives, and policy-gradient learning in machine
translation.

\field{Question:}
Can the sign of a pre-optimization gradient-alignment test predict whether
policy-gradient optimization of a learned machine-translation reward will
improve or reduce independently measured human quality on new reward models
and held-out data?

\field{Why this question matters:}
Policy-gradient methods optimize the reward they are given.  In realistic
applications, that reward is often a learned proxy for the outcome people
actually care about.  An unbiased, low-variance gradient estimator can
therefore work exactly as intended while moving the policy in a harmful
direction.  Machine translation provides a useful test bed because the policy
can choose among real system outputs, learned reward scores are available, and
held-out human judgments can evaluate the selected translations without
entering the learning signal.

\field{Preliminary evidence:}
I have already built the finite-action bandit and compared exact policy-gradient
updates with sampled REINFORCE.  The learned proxy rewards increase, but after
extended optimization all \SeedCellsBelowFinal{} of \SeedCells{} independently
trained learned-reward instances finish below uniform selection in human
quality.  The failure also persists after replacing system-identity logits
with a shared source--candidate-text scorer: under the all-data reward, held-out
human quality changes by \AllTextChange{} points
$[\AllTextLow{},\AllTextHigh{}]$ on unseen documents and systems.  A
retrospective analysis suggests that the initial alignment between the reward
gradient and a human-quality gradient may predict the direction of this
failure.  These are development findings used to form the hypothesis, rather
than prospective confirmation.

\field{MDP:}
This is a contextual bandit, or a one-decision MDP.
\begin{itemize}
    \item \textbf{State/observation:} a source sentence and the candidate
    translations available for that source.
    \item \textbf{Action:} select one candidate translation produced by a
    machine-translation system.
    \item \textbf{Transition:} the episode terminates immediately after the
    selection.
    \item \textbf{Reward:} the score from a learned translation reward model.
    Human ratings are reserved for evaluation and never supplied as reward.
    \item \textbf{Horizon/discount:} a single decision, so discounting is
    irrelevant.
\end{itemize}

\field{Data and feedback available:}
The development study uses released shared-task sources, candidate system
outputs, and human judgments.  The prospective test uses new reward-training
seeds and two separate human-rated candidate datasets.  Within each dataset,
calibration documents supply the small amount of human feedback needed for the
alignment test.  Disjoint documents and systems supply final evaluation.  Human
scores never enter reward optimization, checkpoint selection, or stopping.

\field{Main course connection / assumption under study:}
The main connection is the score-function estimator and REINFORCE.  The key
assumption under study is that the scalar reward is aligned with the desired
outcome.  Because every action can be enumerated in this finite environment, I
can compute the exact expected objective
\[
J_r(\theta\mid x)=\sum_{a}\pi_\theta(a\mid x)r(x,a)
\]
and its exact gradient.  This creates a known reference for checking sampled
REINFORCE updates and separates estimator error from reward misspecification.

\field{Method I will try first:}
Before reading confirmation outcomes, I have frozen a decision rule.  For each
new reward checkpoint, I compute the cosine between its initial policy-gradient
direction and a human-quality gradient estimated from the calibration sample.
If the cosine is positive, the rule optimizes a shared source--candidate-text
policy; otherwise it keeps the initial uniform policy.  Exact finite-action
gradients isolate reward alignment from estimator noise, while the completed
REINFORCE comparison connects the result to sampled policy gradients.  The
same frozen rule is then evaluated on disjoint documents and unseen systems.

\field{Baseline:}
The decision baselines are always optimize, always retain the initial policy,
and a cheap best-of-$N$ screening rule on the same calibration candidate pool.
Uniform selection remains the initial policy, and the human-best candidate is
reported only as an evaluation ceiling.

\field{What I will measure:}
The primary outcome is decision regret: the human-quality loss relative to the
better of keeping the initial policy and optimizing it.  I will compare mean
regret, harmful optimize decisions, and missed beneficial optimizations across
the frozen alignment rule and all baselines.  Supporting plots show held-out
human quality and proxy reward against matched policy displacement.  Uncertainty
is aggregated over reward-training seeds and evaluation documents.

\field{What I expect to teach the class:}
REINFORCE can be mathematically correct and reliably increase its stated reward
while harming the outcome that motivated the reward.  The completed experiments
separate estimator correctness from objective validity.  The prospective test
then asks a practical follow-up: whether a small, pre-optimization alignment
check can tell us when the safest action is to leave the policy alone.

\end{document}
